% Source: https://artofproblemsolving.com/community/c3h1490800_another_mathcounts_problem
% Stable identity: aops-topic-1490800-q01
Put $P=(0,0)$, $A=(-8,0)$, $B=(8,0)$, $D=(0,7)$. The reflection in line $DP$ (the $y$-axis) swaps $A$ and $B$, so it swaps $C$ and $E$: if $C=(x,y)$ then $E=(-x,y)$, with $y>0$ by convexity (all vertices lie on the same side of $AB$).

\textbf{The conditions.} The direction vectors are $\overrightarrow{BC}=(x-8,\,y)$ and $\overrightarrow{AE}=(8-x,\,y)$. Perpendicularity gives
\[
-(x-8)^2+y^2=0,\qquad\text{so } y=|x-8| .
\]
And $BC=4$ gives $(x-8)^2+y^2=16$, so $2(x-8)^2=16$, $|x-8|=2\sqrt2$ and $y=2\sqrt2$. There are two possibilities, $x=8-2\sqrt2$ and $x=8+2\sqrt2$, and both give convex pentagons:

\begin{asy}
usepackage("amsmath");
size(5cm);
defaultpen(fontsize(10pt));
pair ft(pair P,pair U,pair V){return U+dot(P-U,V-U)/dot(V-U,V-U)*(V-U);}
path ra(pair P,pair u,pair v,real s=0.2){return P+s*unit(u)--P+s*unit(u)+s*unit(v)--P+s*unit(v);}
real s=-1; pair A=(-8,0),B=(8,0),P=(0,0),D=(0,7);
pair C=(8+s*2*sqrt(2),2*sqrt(2)), Ee=(-C.x,C.y), Q=(0,-s*8);
if(s>0){ draw(B--Q,gray+dashed); draw(A--Q,gray+dashed); draw(ra(Q,B-Q,A-Q,0.8),gray); }
else { draw(arc(B,1.6,135,180),gray); draw(arc(A,1.6,0,45),gray);
 label("$45^\circ$",B+3.1*dir(160),fontsize(8pt)); label("$45^\circ$",A+3.1*dir(20),fontsize(8pt)); }
draw(D--P,red+dashed); draw(ra(P,B,D,0.7),red);
draw(A--B--C--D--Ee--cycle,linewidth(0.9));
dot(A);dot(B);dot(C);dot(D);dot(Ee);dot(P);
label("$A$",A,SW);label("$B$",B,SE);label("$C$",C,dir(30));label("$D$",D,N);label("$E$",Ee,W);label("$P$",P,S);
label("$4$",(B+C)/2,dir(45)); label("$7$",(0,3.5),E,red);
label("$AB=16$",(0,-2));
label("\mbox{}",(min(currentpicture,user=true)+max(currentpicture,user=true))/2,black); // keeps the text after the figure black when LaTeX typesets the labels inline
\end{asy}

\begin{asy}
usepackage("amsmath");
size(5cm);
defaultpen(fontsize(10pt));
pair ft(pair P,pair U,pair V){return U+dot(P-U,V-U)/dot(V-U,V-U)*(V-U);}
path ra(pair P,pair u,pair v,real s=0.2){return P+s*unit(u)--P+s*unit(u)+s*unit(v)--P+s*unit(v);}
real s=1; pair A=(-8,0),B=(8,0),P=(0,0),D=(0,7);
pair C=(8+s*2*sqrt(2),2*sqrt(2)), Ee=(-C.x,C.y), Q=(0,-s*8);
if(s>0){ draw(B--Q,gray+dashed); draw(A--Q,gray+dashed); draw(ra(Q,B-Q,A-Q,0.8),gray); }
else { draw(arc(B,1.6,135,180),gray); draw(arc(A,1.6,0,45),gray);
 label("$45^\circ$",B+3.1*dir(160),fontsize(8pt)); label("$45^\circ$",A+3.1*dir(20),fontsize(8pt)); }
draw(D--P,red+dashed); draw(ra(P,B,D,0.7),red);
draw(A--B--C--D--Ee--cycle,linewidth(0.9));
dot(A);dot(B);dot(C);dot(D);dot(Ee);dot(P);
label("$A$",A,dir(200));label("$B$",B,dir(-80));label("$C$",C,E);label("$D$",D,N);label("$E$",Ee,W);label("$P$",P,S);
label("$4$",(B+C)/2,dir(-45)); label("$7$",(0,3.5),E,red);
label("$AB=16$",(0,-10));
label("\mbox{}",(min(currentpicture,user=true)+max(currentpicture,user=true))/2,black); // keeps the text after the figure black when LaTeX typesets the labels inline
\end{asy}

In the first, angles $ABC$ and $BAE$ are $45^\circ$, so lines $BC$ and $AE$ meet at a right angle above $D$ (at $(0,8)$); in the second, these angles are $135^\circ$ and the lines meet at a right angle below $AB$ (at $(0,-8)$, grey dashed).

\textbf{The area.} By symmetry the area is twice the area of quadrilateral $PBCD$ with vertices $(0,0)$, $(8,0)$, $(x,y)$, $(0,7)$. By the shoelace formula,
\[
[PBCD]=\dfrac{1}{2}\left(8y+7x\right),\qquad [ABCDE]=8y+7x=16\sqrt2+7x .
\]
\begin{itemize}
\item $x=8-2\sqrt2$: area $16\sqrt2+56-14\sqrt2=56+2\sqrt2$.
\item $x=8+2\sqrt2$: area $16\sqrt2+56+14\sqrt2=56+30\sqrt2$.
\end{itemize}
So the area is $\boxed{56+2\sqrt2}$ (when angle $ABC$ is acute, $45^\circ$) or $\boxed{56+30\sqrt2}$ (when angle $ABC$ is obtuse, $135^\circ$).
